
\begin{frame}{Le \textit{Streaming} avec Flink}

On a affaire à un \alert{stream}, flux de données de taille (conceptuellement) infinie.

\vfill

Les éléments de ce flux sont des \alert{item}s dans la terminologie Flink. Un item est de type quelconque,
on doit parfois lui attribuer une clé de regroupement.

\vfill

Et on applique à ce flux et à ces items des \alert{transformations}, en continu

\vfill
\begin{blocgauche}
Contrairement à Spark, Flink traite nativement les flux: les items sont pris en compte
un par un, instantanémenent (alors que Spark fait du micro-batch, ou pseudo-streaming)
\end{blocgauche}
\end{frame}


\begin{frame}{Les \textit{workflow} de Flink}

La source est typiquement un gestionnaire de flux comme Kafka.
\vfill

\figSlideWithSize{flink_dataflow}{10cm}

\vfill

La cible est typiquement une base de données. 

\vfill
\begin{blocgauche}
Les opérateurs (voir la liste dans support de cours) sont à peu près ceux que l'on trouve dans Pig ou Spark.
\end{blocgauche}
\end{frame}

%%%%%%%%%%%%%
\begin{frame}{Notre simulateur de flux}

Pour expérimenter facilement, nous utilisons un simulateur de flux écrit en Python, \texttt{Simflux.py}. 

\vfill

\begin{redbullet}

  \item Pour engendrer un flux d'entiers sur le port 9000 (par défaut)\,: 
  
   \texttt{python3 SimFlux.py}

\item Pour engendrer un flux de lignes extraites d'un fichier texte\,: 

 \texttt{python3 SimFlux.py {-}-source nom\_fichier}

\item Pour engendrer un flux des contenus de fichier d'un répertoire\,:  

  \texttt{python3 SimFlux.py {-}-source répertoire}

\end{redbullet}

\vfill

On va tester la façon dont Flink consomme ce flux en entrée.
\end{frame}


\begin{frame}[fragile]{Un premier exemple avec l'interpréteur Scala}

Au préalable, lancer notre générateur de flux sur le port 9000.

\vfill

\begin{minted}{scala}
// Déclaration du flux de données. NB: senv est l'environnement de streaming
val stream = senv.socketTextStream("localhost", 9000, '\n')
// À chaque entier en entrée, on ajoute 2 (pourquoi pas)
val w = stream.map ({ x => x.toInt + 2 } )
// Voyons ce que cela donne
w.print()
// Workflow défini: on l'exécute
senv.execute("Mon premier traitement de flux ")
\end{minted}

\vfill

\begin{blocgauche}
Notez bien que l'exécution est ``paresseuse'': on ne déclenche le \textit{workflow} que quand le \textit{sink} est défini.
\end{blocgauche}
\end{frame}

%%%%%%%%%%%%
\begin{frame}[fragile]{Un workflow = séquence de transformations}

S'exprime très simplement en Scala.

\vfill

Ci-dessous\,: on crée un Tuple scala avec l'entier reçu, puis on lui associe son double, puis
on affiche.

\vfill
\begin{minted}{scala}
stream.map ( { x => Tuple1(x.toInt) } )
      .map( {y => (y._1, y._1 * 2) } )
      .print()
\end{minted}

\vfill
\begin{blocgauche}
Astuce: utiliser \texttt{:paste} puis \texttt{CTRL D} pour entrer des commandes multi-lignes.
\end{blocgauche}
\end{frame}



%%%%%%%%%
\begin{frame}[fragile]{Traitons des flux d'objets}

Plus intéressant: on transforme le flux en entrée en un type donné (comme dans Pig)

\vfill

\begin{minted}{scala}
case class MonDouble(leReel: Float, sonDouble: Double)
val stream = senv.socketTextStream("localhost", 9000, '\n')
val w = stream.map ( { x => Tuple1(x.toFloat) } )
                .map( {y => MonDouble(y._1, y._1 * 2) } )
val fluxFiltre = w.filter ({_.leReel > 1000})
fluxFiltre.print()
senv.execute("Mon premier traitement de flux ")
\end{minted}

\vfill
\begin{blocgauche}
Important: \textit{Streaming} = pipelinage. \alert{Il n'y a pas de matérialisation des résultats intermédiaires}.
\end{blocgauche}
\end{frame}


%%%%%%%%%%%%
\begin{frame}[fragile]{Passons à des chaînes de caractères}

Ci-dessous\,: après avoir lancé le simulateur en mode 'lecture de fichier texte',
on traite les lignes de texte constituant le flux.

\vfill
\begin{minted}{scala}
     val stream = senv.socketTextStream("localhost", 9000, '\n')
     val w = stream.flatMap ({ str => str.split("\\W+") }).print()
     senv.execute("Je découpe en mots")
\end{minted}

\vfill

Nous sommes prêts à compter les mots.
\end{frame}

\begin{frame}[fragile]{Le fameux compteur de mots}

Voici la chaîne de traitement pour compter les mots d'un flux en entrée.
\vfill

\begin{minted}{scala}
case class CompteurMot(mot: String, compteur: Int)
 val stream = senv.socketTextStream("localhost", 9000, '\n')
 val w = stream.flatMap ({ str => str.split("\\W+") })
        .map({ CompteurMot(_, 1) })
        .keyBy("mot")
        .sum("compteur")
        .print()
senv.execute("Le compteur de mots")
\end{minted}

\vfill
\begin{blocgauche}
Notez.  (i) l'opérateur \texttt{keyBy} qui regroupe des nuplets Scala sur un champ (ici, le champ ``mot'') 
et crée donc autant de flux que de mots; (ii) le compteur ne fait que s'incrémenter.
\end{blocgauche}
\end{frame}


\begin{frame}[fragile]{Le même, en utilisant \texttt{reduce}}

Que signifie \texttt{reduce} dans un contexte de flux? C'est l'accumulation
successive des valeurs rencontrées, groupées d'après une clé.

\vfill

\begin{minted}{scala}
case class CompteurMot(mot: String, compteur: Int)
val stream = senv.socketTextStream("localhost", 9000, '\n')
val mots = stream.flatMap ({ str => str.split("\\W+") })
                 .map({ CompteurMot(_, 1) }).keyBy("mot")
val compte = mots.reduce( (acc, occ)  
       => {CompteurMot (acc.mot, acc.compteur + 1) }).print()
senv.execute("Le compteur de mots")
\end{minted}

\vfill
\begin{blocgauche}
Et voilà\,! Nous savons traiter des flux \emph{infinis}. Il reste à savoir
les découper en \alert{fenêtres}.
\end{blocgauche}
\end{frame}

